% ================= 第 5 章 (前半): 5.1 - 5.4 =================

\section{主要结果}

我们证明引言的定理 \ref{thm:1.1.4} 和 \ref{thm:1.1.2}.

\subsection{关于 \ref{sec:3.1} 的空间 $\widehat{H}{}^{1}(K_1^{p}(M))$ 和 $\widehat{H}{}_{c}^{1}(K_1^{p}(M))$ 的预备结果}

设 $M \ge 1$ 为与 $p$ 互素的整数, 对 $r \ge 1$, $Y_1(Mp^{r})$ 为与同余子群 $\Gamma_1(Mp^{r})$ 对应的开模曲线.
记:
\[
\widehat{H}_c := \varprojlim_n \frac{\varinjlim_r H_{c}^{1}(Y_1(Mp^{r}),\Zp)}{p^{n}\,\varinjlim_r H_{c}^{1}(Y_1(Mp^{r}),\Zp)}.
\]
它自然地是 Hecke $\Zp$-模, 还配备算子 $U_p$ 的作用, 且 $\widehat{H}_c \otimes_{\Zp} \Qp$ 是 Banach 空间.
投影 $Y_1(M;p^{r}) \to Y_1(Mp^{r})$ 诱导连续 $\mathbb{T}$-等变映射:
\begin{equation}
\label{eq:23}
\widehat{H}_c \to \widehat{H}_{c}^{1}(K_1^{p}(M))^{P(\Zp)},
\end{equation}
其中 $P(\Zp) := \left\{\begin{pmatrix} a & b \\ 0 & 1 \end{pmatrix},\ a \in \Zp^{\times},\ b \in \Zp\right\}$ 是引言中引入的子群 $P$ 的 $\Zp$-点群, 它把 $\widehat{H}_c$ 上 $U_p$ 的作用与 $\widehat{H}_{c}^{1}(K_1^{p}(M))^{P(\Zp)}$ 上 $p\pi_{\wp}$ 的作用交织在一起, 其中 $\wp := \begin{pmatrix} p & 0 \\ 0 & 1 \end{pmatrix}$.
对一切 $r \ge 1$, 记 $G_r$ 为约化 $\GL_2(\Zp) \to \GL_2(\Zp/p^{r}\Zp)$ 的核.

\begin{proposition}\label{prop:5.1.1}
映射 \eqref{eq:23} 是同构.
\end{proposition}

\begin{proof}
论证类似于 \cite[Th.~2.1.12]{ref:24} 的证明.
取 $r_0$ 足够大使得 $\Gamma_1(Mp^{r_0})$ 无挠, 设 $X_{r_0}$ 为 $Y_1(Mp^{r_0})$ 的维数 $1$ 单纯子复形, 它是 $Y_1(Mp^{r_0})$ 的形变收缩核.
设:
\[
G_{r_0}^{1} := \left\{ \begin{pmatrix} a & b \\ c & d \end{pmatrix} \in \GL_2(\Zp) \;\middle|\; (c,d) \equiv (0,1) \bmod p^{r_0} \right\},
\]
则对 $r > r_0$, $Y_1(M;p^{r})$ 是 $Y_1(Mp^{r_0})$ 的 Galois 覆盖, Galois 群为 $G_{r_0}^{1}/G_r$ (右作用).
设 $X_r$ 为 $X_{r_0}$ 在 $Y_1(M;p^{r})$ 中的逆像, 它仍是 $X_{r_0}$ 的 Galois 覆盖, Galois 群为 $G_{r_0}^{1}/G_r$.
记 $T_{\bullet}(r_0)$ 为 $X_{r_0}$ 的三角剖分 (其中 $T_i(r_0)$ 是次数 $i$ 的单纯形的集合, $i \in \{0,1\}$), $T_{\bullet}(r)$ 为 $X_r$ 的逆像三角剖分 $T_{\bullet}(r_0)$.
$G_{r_0}^{1}/G_r$ 在 $X_{r_0}$ 之上 $X_r$ 上的右作用诱导其在 $T_{\bullet}(r_0)$ 上的右作用.
对 $r \ge r_0$, 设 $S_{\bullet}(r,\Zp/p^{n}\Zp)$ 为与三角剖分 $T_{\bullet}(r)$ 相关联的系数在 $\Zp/p^{n}\Zp$ 中的链复形: 它是 $(\Zp/p^{n}\Zp)[G_{r_0}^{1}/G_r]$-模复形.
当 $r$ 变化时, 有链复形的归纳系 $\bigl(S_{\bullet}(r,\Zp/p^{n}\Zp)\bigr)_r$, 记:
\[
S_{\bullet}(\Zp/p^{n}\Zp) := \varinjlim_r S_{\bullet}(r,\Zp/p^{n}\Zp).
\]
对 $n$ 取射影极限, 得到 $\Zp[[G_{r_0}^{1}]]$-模复形:
\[
\widetilde{S}_{\bullet} := \varprojlim_n S_{\bullet}(\Zp/p^{n}\Zp).
\]
类似地, 设 $X_{1,r}$ 为 $X_{r_0}$ 在 $Y_1(Mp^{r})$ 中的逆像, $T_{\bullet}(1,r)$ 为 $X_{1,r}$ 的逆像三角剖分 $T_{\bullet}(r_0)$.
对 $r \ge r_0$, 设 $R_{\bullet}(r,\Zp/p^{n}\Zp)$ 为与三角剖分 $T_{\bullet}(1,r)$ 相关联的系数在 $\Zp/p^{n}\Zp$ 中的链复形, 并设 $R_{\bullet}(\Zp/p^{n}\Zp) := \varinjlim_r R_{\bullet}(r,\Zp/p^{n}\Zp)$, $\widetilde{R}_{\bullet} := \varprojlim_n R_{\bullet}(\Zp/p^{n}\Zp)$.
由于三角剖分 $T_{\bullet}(r)$ 只涉及 $0$ 维和 $1$ 维单纯形, 复形 $S_{\bullet}(\Zp/p^{n}\Zp)$ 只在次数 $0$ 和 $1$ 有项, 故有正合列:
\[
0 \to H_1(S_{\bullet}(\Zp/p^{n}\Zp)) \to S_1(\Zp/p^{n}\Zp) \to S_0(\Zp/p^{n}\Zp).
\]
射影极限函子左正合, 由此得到正合列:
\[
0 \to \varprojlim_n H_1(S_{\bullet}(\Zp/p^{n}\Zp)) \to \widetilde{S}_1 \to \widetilde{S}_0,
\]
从而有同构:
\begin{equation}
\label{eq:24}
H_1(\widetilde{S}_{\bullet}) \xrightarrow{\sim} \varprojlim_n H_1(S_{\bullet}(\Zp/p^{n}\Zp)).
\end{equation}
归纳极限函子正合, 有:
\begin{equation}
\label{eq:25}
H_1(S_{\bullet}(\Zp/p^{n}\Zp)) \xrightarrow{\sim} \varinjlim_r H_1(Y_1(M;p^{r}),\Zp/p^{n}\Zp) \xrightarrow{\sim} \Bigl(\varinjlim_r H_1(Y_1(M;p^{r}),\Zp)\Bigr) \otimes_{\Zp} \Zp/p^{n}\Zp.
\end{equation}
结合 \eqref{eq:24} 和 \eqref{eq:25}, 推出 $H_1(\widetilde{S}_{\bullet})$ 同构于 $\varinjlim_r H_1(Y_1(M;p^{r}),\Zp)$ 的 $p$ 进完备化, 从而由 Poincar\'e 对偶同构于 $\widehat{H}_{c}^{1}(K_1^{p}(M))$.
若让 $g \in G_{r_0}^{1}$ 经由 $g^{-1}$ 的右作用左作用在 $H_1(\widetilde{S}_{\bullet})$ 上, 这实际上是 $G_{r_0}^{1}$-模同构.

通过与上一段类似的推理, 得到 $H_1(\widetilde{R}_{\bullet})$ 同构于 $\widehat{H}_c$.
另一方面, 由复形本身显然有同构 $\widetilde{R}_{\bullet} \xrightarrow{\sim} (\widetilde{S}_{\bullet})^{P(\Zp)}$, 因此由于 $P(\Zp)$-不变量函子左正合, 对正合列 $0 \to H_1(\widetilde{S}_{\bullet}) \to \widetilde{S}_1 \to \widetilde{S}_0$ 取 $P(\Zp)$-不变量得到同构 $H_1(\widetilde{R}_{\bullet}) \xrightarrow{\sim} H_1(\widetilde{S}_{\bullet})^{P(\Zp)}$.
结合本段与上一段的结论, 得到所求同构 $\widehat{H}_c \xrightarrow{\sim} \widehat{H}_{c}^{1}(K_1^{p}(M))^{P(\Zp)}$.
\end{proof}

现在我们研究扩张空间 $\Ext_{\mathfrak{b}(\Qp)}^{1}\bigl(\psi,\widehat{H}{}^{1}(K_1^{p}(M))_{L,\mathrm{an}}\bigr)$ 和 $\Ext_{\mathfrak{b}(\Qp)}^{1}\bigl(\psi,\widehat{H}{}_{c}^{1}(K_1^{p}(M))_{L,\mathrm{an}}\bigr)$, 其中 $\psi$ 是 $\mathfrak{b}(\Qp)$ 的取值于 $\Qp$ 的有限扩张 $L$ 的特征标, $\widehat{H}{}^{1}(K_1^{p}(M))_{L,\mathrm{an}}$ 和 $\widehat{H}{}_{c}^{1}(K_1^{p}(M))_{L,\mathrm{an}}$ 表示可容许 $\GL_2(\Qp)$-Banach 空间 $\widehat{H}{}^{1}(K_1^{p}(M))_L$ 和 $\widehat{H}{}_{c}^{1}(K_1^{p}(M))_L$ 的局部解析向量 (参见 \cite[\S7]{ref:45}).

\begin{proposition}\label{prop:5.1.2}
对上述任何特征标 $\psi$, 有:
\[
\Ext_{\mathfrak{b}(\Qp)}^{1}\bigl(\psi,\widehat{H}{}^{1}(K_1^{p}(M))_{L,\mathrm{an}}\bigr) = 0.
\]
\end{proposition}

\begin{proof}
论证基于 \cite[Th.~2.1.5]{ref:24} 证明中的构造.
取 $r_0$ 使得 $\Gamma_1(M) \cap \Gamma(p^{r_0})$ 无挠, 设 $X_{r_0}$ 为 $Y_1(M;p^{r_0})$ 的维数 $1$ 单纯子复形, 是 $Y_1(M;p^{r_0})$ 的形变收缩核.
若 $r \ge r_0$, $Y_1(M;p^{r})$ 是 $Y_1(M;p^{r_0})$ 的 Galois 覆盖, Galois 群为 $G_{r_0}/G_r$ (右作用).
设 $X_r$ 为 $X_{r_0}$ 在 $Y_1(M;p^{r})$ 中的逆像, 则 $X_r$ 是 $X_{r_0}$ 的 Galois 覆盖, Galois 群为 $G_{r_0}/G_r$.
然后完全像引理 \ref{prop:5.1.1} 的证明那样定义三角剖分 $T_{\bullet}(r_0)$ 和 $T_{\bullet}(r)$.
对 $r \ge r_0$, 设 $S^{\bullet}(r,\Zp/p^{n}\Zp)$ 为与 $T_{\bullet}(r)$ 相关联的系数在 $\Zp/p^{n}\Zp$ 中的上链复形.
与前一证明类似地定义复形 $S^{\bullet}(\Zp/p^{n}\Zp)$ 和 $\widetilde{S}^{\bullet}$.
特别地, 后者是 $\Zp[[G_{r_0}]]$-模复形.
由 \cite[Th.~2.1.5]{ref:24}, 每个空间 $\widetilde{S}^{i}$ 都同构于某个整数 $s_i \ge 0$ 的 $G_{r_0}$-模 $\mathcal{C}^{0}(G_{r_0},\Zp)^{s_i}$, 且有同构 $H^{i}(\widetilde{S}^{\bullet}) \xrightarrow{\sim} \widehat{H}{}^{i}(K_1^{p}(M))$ ($i \in \{0,1\}$).
由于"局部解析向量"函子正合 (\cite[Th.~7.1]{ref:45}), 得到正合列:
\[
0 \to \widehat{H}{}^{0}(K_1^{p}(M))_{L,\mathrm{an}} \to \mathcal{C}^{\mathrm{an}}(G_{r_0},L)^{s_0} \to \mathcal{C}^{\mathrm{an}}(G_{r_0},L)^{s_1} \to \widehat{H}{}^{1}(K_1^{p}(M))_{L,\mathrm{an}} \to 0.
\]
由于中间两项是 $\mathfrak{b}(\Qp)$-零调的且关于特征标 $-\psi$ 的扭曲不变, 简单的图追踪给出同构:
\begin{equation}
\label{eq:26}
\Ext_{\mathfrak{b}(\Qp)}^{1}\bigl(\psi,\widehat{H}{}^{1}(K_1^{p}(M))_{L,\mathrm{an}}\bigr)
= H^{1}\bigl(\mathfrak{b}(\Qp),\widehat{H}{}^{1}(K_1^{p}(M))_{L,\mathrm{an}}(-\psi)\bigr)
\xrightarrow{\sim} H^{3}\bigl(\mathfrak{b}(\Qp),\widehat{H}{}^{0}(K_1^{p}(M))_{L,\mathrm{an}}(-\psi)\bigr).
\end{equation}
(其中"$(-\psi)$"表示用 $-\psi$ 扭曲.)
但 $\widehat{H}{}^{0}(K_1^{p}(M))_{L,\mathrm{an}} \xrightarrow{\sim} \mathcal{C}^{\mathrm{an}}(\Zp^{\times},L)$, 故 $H^{1}\bigl(\mathfrak{z}(\Qp),\widehat{H}{}^{0}(K_1^{p}(M))_{L,\mathrm{an}}(-\psi)\bigr) = 0$.
另一方面, $\mathfrak{b}(\Qp)$ 是 $\mathfrak{z}(\Qp)$ 与一个维数 $2$ 的 Lie 代数 $\mathfrak{l}(\Qp)$ 的直和 (后者由形如 $\begin{pmatrix} * & * \\ 0 & 0 \end{pmatrix}$ 的矩阵构成的 Lie 代数).
回顾维数 $d$ 的 Lie 代数的上同调在次数 $> d$ 处为零.
由于 $\mathfrak{z}(\Qp)$ 维数 $1$, $\mathfrak{l}(\Qp)$ 维数 $2$, Lie 代数上同调的 K\"unneth 公式推出: 对任何 $\mathfrak{b}(\Qp)$-模 $M$, 有 $H^{3}(\mathfrak{b}(\Qp),M) \simeq H^{1}(\mathfrak{z}(\Qp),M) \otimes H^{2}(\mathfrak{l}(\Qp),M)$ (唯一的可能非零项).
取 $M = \widehat{H}{}^{0}(K_1^{p}(M))_{L,\mathrm{an}}(-\psi)$, 由于 $H^{1}(\mathfrak{z}(\Qp),M) = 0$, 确实得到 $H^{3}(\mathfrak{b}(\Qp),M) = 0$, 从而由 \eqref{eq:26} 得到结果.
\end{proof}

\begin{corollary}\label{cor:5.1.3}
对上述任何特征标 $\psi$, Hecke 模 $\Ext_{\mathfrak{b}(\Qp)}^{1}\bigl(\psi,\widehat{H}{}_{c}^{1}(K_1^{p}(M))_{L,\mathrm{an}}\bigr)$ 是 Eisenstein 的.
\end{corollary}

\begin{proof}
由"局部解析向量"函子的正合性 (\cite[Th.~7.1]{ref:45}), \eqref{eq:7} 的满射性和命题 \ref{prop:5.1.2}, 可见 Hecke 模 $\Ext_{\mathfrak{b}(\Qp)}^{1}\bigl(\psi,\widehat{H}{}_{c}^{1}(K_1^{p}(M))_{L,\mathrm{an}}\bigr)$ 是 $\Ext_{\mathfrak{b}(\Qp)}^{1}\bigl(\psi,\widehat{M}_{L,\mathrm{an}}\bigr)$ 的商.
推论则从推论 \ref{cor:3.1.3} 证明中的同一论证并利用引理 \ref{lem:3.1.2} 得出.
\end{proof}

\subsection{关于 \cite{ref:22} 和 \cite{ref:25} 的 Jacquet 函子以及 \cite{ref:24} 关于 $\Q$ 上 $\GL_2$ 的结果的一些回顾}

设 $V$ 是局部凸拓扑 $L$-向量空间 (其中 $L$ 是 $\Qp$ 的有限扩张), 配备 $B(\Qp)$ 的局部解析作用.
在 \cite[Def.~3.4.5]{ref:22} 中定义了一个紧型局部凸拓扑 $L$-向量空间 $J_{B(\Qp)}(V)$, 配备 $T(\Qp)$ 的局部解析表示, 称为 $V$ 的 Jacquet 模.
这里不回顾它的定义, 而只需注意: 若 $\chi_1 \otimes \chi_2 \in \widehat{T}(L)$ (\ref{sec:2.1}) 且 $J_{B(\Qp)}^{\chi_1 \otimes \chi_2}(V)$ 是 $J_{B(\Qp)}(V)$ 中 $T(\Qp)$ 经 $\chi_1 \otimes \chi_2$ 作用的子空间, 则有自然同构 (参见 \cite[Prop.~3.4.9]{ref:22}):
\begin{equation}
\label{eq:27}
J_{B(\Qp)}^{\chi_1 \otimes \chi_2}(V) \xrightarrow{\sim} V^{N(\Zp),T(\Qp)^{+} = \chi_1 \otimes \chi_2},
\end{equation}
其中 $V^{N(\Zp),T(\Qp)^{+} = \chi_1 \otimes \chi_2}$ 是 $V^{N(\Zp)}$ 中对一切 $t \in T(\Qp)^{+}$, $\pi_t$ 按 $(\chi_1 \otimes \chi_2)(t)$ 倍乘作用的子空间.

Jacquet 模 $J_{B(\Qp)}\bigl(\widehat{H}{}_{c}^{1}(K_1^{p}(M))^{\mathrm{an}}\bigr)$ 是 \cite{ref:21} 意义下 $T(\Qp)$ 的本质可容许局部解析表示, 因此对应 $\widehat{T}$ 上一个凝聚层, 记为 $\mathcal{E}_c$ (参见 \cite[\S4.4]{ref:24}).
Hecke 算子 $T_{\ell}$ 和 $S_{\ell}^{\pm 1}$ 在 $\widehat{H}{}_{c}^{1}(K_1^{p}(M))$ 上 (对 $\ell \nmid Mp$) 的作用诱导 $\mathcal{E}_c$ 上的作用.
如 \cite[\S4.4]{ref:24}, 记 $\mathcal{A}_c \subset \End(\mathcal{E}_c)$ 为由 $T_{\ell}$ 和 $S_{\ell}^{\pm 1}$ 生成的凝聚子代数额层, $\Spec \mathcal{A}_c$ 为 $\mathcal{A}_c$ 在 $\widehat{T}$ 上的相对 Spec.
因此 $\Spec \mathcal{A}_c$ 是刚性解析空间, 配备连续注入:
\[
\Spec \mathcal{A}_c \hookrightarrow \widehat{T} \times \Spec \Qp\bigl[\{T_{\ell},S_{\ell}^{\pm 1}\}_{\ell \nmid Mp}\bigr]
\simeq (\mathcal{W} \times \mathbb{G}_m)^{2} \times \Spec \Qp\bigl[\{T_{\ell},S_{\ell}^{\pm 1}\}_{\ell \nmid Mp}\bigr].
\]
我们通常用 $(\chi_2|\,| \otimes \chi_1|\,|^{-1},\lambda)$ 记 $\Spec \mathcal{A}_c$ 的一个点, 其中 $\chi_2|\,| \otimes \chi_1|\,|^{-1}$ 表示 $\widehat{T}$ 的相应点, $\lambda$ 为相应的 Hecke 特征值系 (用 $|\,| \otimes |\,|^{-1}$ 扭曲的选择将在 \ref{sec:5.5} 解释).
空间 $\Spec \mathcal{A}_c$ 关于 $\mathcal{W}$ 的元素是"对角"扭曲不变的 (参见 \cite{ref:24}, (4.4.8)): 点 $(\chi_2|\,| \otimes \chi_1|\,|^{-1},\lambda)$ 被 $\chi \in \mathcal{W}$ 扭曲后是点 $(\chi\chi_2|\,| \otimes \chi\chi_1|\,|^{-1},\chi\lambda)$.

另一方面由 \eqref{eq:27} 推出:

\begin{proposition}\label{prop:5.2.1}
有 Hecke-等变同构:
\[
J_{B(\Qp)}^{\chi_2|\,| \otimes \chi_1|\,|^{-1}}\bigl(\widehat{H}{}_{c}^{1}(K_1^{p}(M))_{L,\mathrm{an}}^{\lambda}\bigr)
\xrightarrow{\sim} \widehat{H}{}_{c}^{1}(K_1^{p}(M))_{L,\mathrm{an}}^{N(\Zp),T(\Qp)^{+} = \chi_2|\,| \otimes \chi_1|\,|^{-1},\lambda}.
\]
\end{proposition}

由于 $\GL_2(\Qp)$ 在 $\widehat{H}{}_{c}^{1}(K_1^{p}(M))$ 上的作用是酉的, 注意对 $\Spec \mathcal{A}_c$ 的任何点 $(\chi_2|\,| \otimes \chi_1|\,|^{-1},\lambda)$ 有 (参见 \cite[Lem.~4.4.2]{ref:22}):
\begin{equation}
\label{eq:28}
|\chi_2(p)| \le 1 \quad \text{且} \quad |\chi_1(p)| = |\chi_2(p)|^{-1}.
\end{equation}

\subsection{刚性空间 $\Spec \mathcal{A}_c$ 与 Coleman 和 Mazur 的 Hecke 曲线 (或曲面) 之间的联系}\label{sec:5.3}

设 $f$ 是如 \ref{sec:3.2} 的模形式 (即权 $k \ge 2$、水平 $N = p^{r}M$ ($r \ge 0$ 为整数)、定义在 $\Qp$ 中 $\Qp$ 的有限扩张 $L$ 上的归一化尖点形式), 具有水平外特征值系 $\lambda$, 且 $(\Sym^{k-2}L^{2})^{\vee} \otimes \pi_p(f) = \pind{\mathrm{lp},\le k-2}_{B(\Qp)}^{\GL_2(\Qp)} \chi_1 \otimes \chi_2$.
由 \eqref{eq:14} 和经典 Jacquet 函子与局部解析 Jacquet 函子的相容性 (\cite[Prop.~4.3.6]{ref:22}), 可见点 $(\chi_2|\,| \otimes \chi_1|\,|^{-1},\lambda)$ 和 $(\chi_1z^{k-2} \otimes \chi_2z^{2-k},\lambda)$ 都在 $\Spec \mathcal{A}_c$ 中.
也可以从 $\pi_p(f)$ 是特殊级数的形式 $f$ 以及从 Eisenstein 级数构造 $\Spec \mathcal{A}_c$ 中的点.
称 $\Spec \mathcal{A}_c$ 中由前述点经 $\mathcal{W}$ 中任意特征标扭曲得到的点为经典点 (注意这与 \cite{ref:21} 的定义略有不同, 那里只允许用 $\Zp^{\times}$ 的局部代数特征标扭曲.
这在下文中不会造成后果, 因为 $\Zp^{\times}$ 的局部代数特征标在 $\mathcal{W}$ 中 Zariski 稠密).
$\Spec \mathcal{A}_c$ 中经典点的 Zariski 闭包 (在刚性解析意义下) 按定义是水平 $M$ 的 Hecke 曲面 (参见 \cite[Prop.~4.4.15]{ref:24}).
它是维数 $2$ 的等维刚性解析簇, 在 $\mathcal{W}$ 的元素作用下稳定.
Hecke 曲面的点 $(\chi_2|\,| \otimes \chi_1|\,|^{-1},\lambda)$ 对应于水平 $M$ 的 $p$ 进模形式 $f$, 它们是有限斜率、温和水平 $M$ 且为 Hecke 特征向量的超收敛模形式经 $\mathcal{W}$ 中元素扭曲所得.
更精确地, 若 $f$ 是有限斜率、整权 $k$、水平 $N = Mp^{r}$ (对 $(M,p) = 1$ 和 $r \ge 0$)、定义在 $\Qp$ 的有限扩张 $L$ 上 (即其 $q$-展开位于 $L \otimes_{\mathcal{O}_L} \mathcal{O}_L[[q]]$ 中)、特征标 $\chi : (\mathbb{Z}/N\mathbb{Z})^{\times} \to L^{\times}$、是 $U_p$ 的特征向量 (特征值为 $\alpha$) 且是 $T_{\ell}$, $S_{\ell}$ 的特征向量 (特征值系为 $\lambda$) 的超收敛形式, 则 $f$ 对应 $(\Spec \mathcal{A}_c)(L)$ 的点:
\[
\bigl(\nr(\alpha/p) \otimes \nr(p/\alpha)\chi_p^{-1}\varepsilon^{2-k},\ \lambda\bigr).
\]
例如, 若 $f$ 是如 \ref{sec:4.3} 对 $\Gamma_1(M)$ 为新的经典形式, $\lambda$ 为相应的 Hecke 特征值系, 则对应 $f(z) - \alpha_p f(pz)$ 的点是 $(\chi_1z^{k-2} \otimes \chi_2z^{2-k},\lambda)$, 对应 $f(z) - \beta_p f(pz)$ 的点是 $(\chi_2|\,| \otimes \chi_1|\,|^{-1},\lambda)$.
同样, 若 $f$ 对 $\Gamma_1(Mp^{r})$ ($r > 0$) 为新且如 \ref{sec:4.4} 在 $p$ 处零斜率, 则对应 $f$ 的点是 $(\chi_1z^{k-2} \otimes \chi_2z^{2-k},\lambda)$, 对应 $f\,|\,w_{p^{r}}$ 的点是 $(\chi_2|\,| \otimes \chi_1|\,|^{-1},\lambda)$ 被 $\chi_0$ 扭曲所得.

Hecke 曲面实际上是 $\Spec \mathcal{A}_c$ 的连通分支之并, 其在 $\Spec \mathcal{A}_c$ 中的补集是维数 $1$ 的不可约分支的不交并, 其中每个分支都只是 $\mathcal{W}$ 作用的一个轨道 (参见 \cite[Prop.~4.4.14]{ref:24} 证明之后的讨论, 关键在于任何维数 $2$ 的分支都包含经典点的 Zariski 稠密集合).
事实上, 我们期望 $\Spec \mathcal{A}_c$ 与 Hecke 曲面重合, 至少在 Eisenstein 轨迹之外如此 (参见 \cite[Thm.~7.5.8]{ref:23}).

称 $\Spec \mathcal{A}_c$ 的点 $(\chi_2|\,| \otimes \chi_1|\,|^{-1},\lambda)$ 为经典权的, 若按定义 \ref{def:2.1.1} 的意义 $\chi_1 \otimes \chi_2$ 是经典权的.
$\Spec \mathcal{A}_c$ 的每个经典点都是经典权的, 但反过来不真 (例如, 不是经典模形式的有限斜率整权超收敛形式给出经典权但非经典的点).

\begin{definition}\label{def:5.3.1}
称 $\Spec \mathcal{A}_c$ 的点 $(\chi_2|\,| \otimes \chi_1|\,|^{-1},\lambda)$ 为序的, 若特征标 $\chi_1$ 和 $\chi_2$ 都是酉的 (即整值).
\end{definition}

\begin{proposition}\label{prop:5.3.2}
$\Spec \mathcal{A}_c$ 的每个序点 $(\chi_2|\,| \otimes \chi_1|\,|^{-1},\lambda)$ 都位于 Hecke 曲面上.
\end{proposition}

\begin{proof}
必要时用 $\chi_2^{-1}|_{\Zp^{\times}}$ 扭曲, 可以假设 $\chi_2$ 非分歧.
则 $J_{B(\Qp)}^{\chi_2|\,| \otimes \chi_1|\,|^{-1}}\bigl(\widehat{H}{}^{1}(K_1^{p}(M))^{\mathrm{an}}_{\lambda}\bigr)$ 的任何元素都由命题 \ref{prop:5.2.1} 特别地对应 $\widehat{H}{}^{1}(K_1^{p}(M))^{P(\Zp)}$ 的一个元素, 从而由命题 \ref{prop:5.1.1} 对应 $\widehat{H}_c$ 的一个元素.
但由 \cite{ref:34}, $\widehat{H}_c$ 的序部分中的任何 Hecke 特征值系都来自一个 $p$ 进模形式.
然后在 \cite{ref:32} 中证明了 (另见对 $p = 2$ 的 \cite{ref:11}) 任何序 $p$ 进模形式都是超收敛的 (与定理 \ref{thm:4.4.8} 证明结尾相同的原理).
\end{proof}

\subsection{Hecke 曲面上"坏"点的定义与研究}\label{sec:5.4}

\begin{definition}\label{def:5.4.1}
$\Spec \mathcal{A}_c$ 的点 $(\chi_2|\,| \otimes \chi_1|\,|^{-1},\lambda)$ 是坏的, 若它是经典权 $k$ 的, 且点 $(\chi_2|\,|z^{1-k} \otimes \chi_1|\,|^{-1}z^{k-1},\lambda)$ 也在 $\Spec \mathcal{A}_c$ 中.
\end{definition}

由命题 \ref{prop:5.2.1} 等价地, $(\chi_2|\,| \otimes \chi_1|\,|^{-1},\lambda)$ 是坏的, 若它是经典权 $k$ 的, 且空间:
\[
\widehat{H}{}_{c}^{1}(K_1^{p}(M))_{L,\mathrm{an}}^{N(\Zp),T(\Qp)^{+} = \chi_2|\,|z^{1-k} \otimes \chi_1|\,|^{-1}z^{k-1},\lambda}
\]
非零 (其中 $L$ 是该点定义于其上的 $\Qp$ 的有限扩张).
定义 \ref{def:5.4.1} 是把 \cite[Def.~4.5.5]{ref:24} 中定义的 Hecke 曲线上的坏点概念推广到 $\Spec \mathcal{A}_c$ 的点.
注意 $\Spec \mathcal{A}_c$ 上坏点的集合在 $\mathcal{W}$ 作用下不变.

\begin{lemma}\label{lem:5.4.2}
若 $(\chi_2|\,| \otimes \chi_1|\,|^{-1},\lambda) \in (\Spec \mathcal{A}_c)(L)$ 是经典权的, 且 $\widehat{H}{}_{c}^{1}(K_1^{p}(M))_{L,\mathrm{an}}^{N(\Zp),T(\Qp)^{+} = \chi_2|\,| \otimes \chi_1|\,|^{-1},\lambda}$ 中存在在 $\SL_2(\Qp)$ 作用下不是局部代数的元素, 则 $(\chi_2|\,| \otimes \chi_1|\,|^{-1},\lambda)$ 是 $\Spec \mathcal{A}_c$ 的坏点.
\end{lemma}

\begin{proof}
假设 $\chi_2|\,| \otimes \chi_1|\,|^{-1}$ 是经典权 $k$ 的.
若 $v \in \widehat{H}{}_{c}^{1}(K_1^{p}(M))_{L,\mathrm{an}}^{N(\Zp),T(\Qp)^{+} = \chi_2|\,| \otimes \chi_1|\,|^{-1},\lambda}$ 不是局部代数的, 则 $X_{-}^{k+1}v$ 非零.
由 \cite[Prop.~4.4.4]{ref:22}, 它属于:
\[
\widehat{H}{}_{c}^{1}(K_1^{p}(M))_{L,\mathrm{an}}^{N(\Zp),T(\Qp)^{+} = \chi_2|\,|z^{1-k} \otimes \chi_1|\,|^{-1}z^{k-1},\lambda},
\]
故 $(\chi_2|\,| \otimes \chi_1|\,|^{-1},\lambda)$ 是坏点.
\end{proof}

\begin{proposition}\label{prop:5.4.3}
设 $(\chi_2|\,| \otimes \chi_1|\,|^{-1},\lambda) \in (\Spec \mathcal{A}_c)(L)$ 是经典权 $k$ 的.
若 $(\chi_2|\,| \otimes \chi_1|\,|^{-1},\lambda)$ 不是经典点, 则它是坏的.
反过来, 若 $(\chi_2|\,| \otimes \chi_1|\,|^{-1},\lambda)$ 是坏的, 则它要么不是经典点, 要么 $|\chi_2(p)| = |p^{k-1}|$.
\end{proposition}

\begin{proof}
由于 $\widehat{H}{}_{c}^{1}(K_1^{p}(M))_{L,\mathrm{an}}$ 中任何 $\SL_2(\Qp)$-局部代数向量都是对某个 $k \ge 2$ 的 \eqref{eq:12} 的像的元素扭曲所得, 可见若 $(\chi_2|\,| \otimes \chi_1|\,|^{-1},\lambda)$ 不是经典点, 则 $\widehat{H}{}_{c}^{1}(K_1^{p}(M))_{L,\mathrm{an}}^{N(\Zp),T(\Qp)^{+} = \chi_2|\,| \otimes \chi_1|\,|^{-1},\lambda}$ 中任何向量都不可能是 $\SL_2(\Qp)$-局部代数的.
因此由引理 \ref{lem:5.4.2}, $(\chi_2|\,| \otimes \chi_1|\,|^{-1},\lambda)$ 是坏的.
反之, 若 $(\chi_2|\,| \otimes \chi_1|\,|^{-1},\lambda)$ 是经典的, 则已知 $|p^{k-1}| \le |\chi_2(p)| \le 1$ (对任何经典点成立).
若第一个不等式是严格的, 则对 $(\chi_2|\,|z^{1-k} \otimes \chi_1|\,|^{-1}z^{k-1},\lambda)$ 应用 \eqref{eq:28} 推出:
\[
\widehat{H}{}_{c}^{1}(K_1^{p}(M))_{L,\mathrm{an}}^{N(\Zp),T(\Qp)^{+} = \chi_2|\,|z^{1-k} \otimes \chi_1|\,|^{-1}z^{k-1},\lambda} = 0,
\]
故 $(\chi_2|\,| \otimes \chi_1|\,|^{-1},\lambda)$ 不是坏的.
\end{proof}

借助命题 \ref{prop:5.4.3}, 可以给出一个识别 Hecke 曲面上坏点的简单而实用的判据.

\begin{proposition}\label{prop:5.4.4}
若 $f$ 是整权 $k \ge 2$、水平 $N = Mp^{r}$ ($r \ge 0$) 的非零超收敛形式, 是 $U_p$, $\mathbb{T}$ 以及 $(\mathbb{Z}/N\mathbb{Z})^{\times}$ 作用的特征向量, 则 Hecke 曲面上对应的点是 $\Spec \mathcal{A}_c$ 的坏点, 当且仅当存在权 $2-k$、水平 $N$、是 $U_p$, $\mathbb{T}$ 以及 $(\mathbb{Z}/N\mathbb{Z})^{\times}$ 作用的特征向量的非零超收敛形式 $g$, 使得 $f$ 与 $\theta^{k-1}(g)$ 有相同的 $U_p$ 特征值、相同的 Hecke 特征值系和相同的 $(\mathbb{Z}/N\mathbb{Z})^{\times}$ 特征标 (关于 $\theta^{k-1}$ 参见 \ref{sec:4}).
\end{proposition}

\begin{proof}
用 \ref{sec:5.3} 的记号, 回顾 $f$ 对应 $\Spec \mathcal{A}_c$ 的点 $\bigl(\nr(\alpha/p) \otimes \nr(p/\alpha)\chi_p^{-1}\varepsilon^{2-k},\lambda\bigr)$.
先假设存在这样的形式 $g$.
容易看出 $g$ 对应 $\Spec \mathcal{A}_c$ 的点 $\bigl(\nr(\alpha/p^{k}) \otimes \nr(p^{k}/\alpha)\chi_p^{-1}\varepsilon^{k},\varepsilon^{k-1}\lambda\bigr)$.
用 $\varepsilon^{1-k} \in \mathcal{W}$ 扭曲, 可见点 $\bigl(\nr(\alpha/p)z^{1-k} \otimes \nr(p/\alpha)\chi_p^{-1}\varepsilon^{2-k}z^{k-1},\lambda\bigr)$ 也在 $\Spec \mathcal{A}_c$ 中, 这表明与 $f$ 关联的点是坏的.
反过来, 假设与 $f$ 关联的点是坏的, 即 $\bigl(\nr(\alpha/p)z^{1-k} \otimes \nr(p/\alpha)\chi_p^{-1}\varepsilon^{2-k}z^{k-1},\lambda\bigr)$ 也在 $\Spec \mathcal{A}_c$ 中.
这次用 $\varepsilon^{k-1}$ 扭曲, 可见 $\bigl(\nr(\alpha/p^{k}) \otimes \nr(p^{k}/\alpha)\chi_p^{-1}\varepsilon^{k},\varepsilon^{k-1}\lambda\bigr)$ 是 $\Spec \mathcal{A}_c$ 的点.
若 $f$ 不是经典的, 由 \cite{ref:13} 和 \cite{ref:14} 已知 $f = \theta^{k-1}(g)$.
因此可以假设 $f$ 是经典的.
对 $f$ 应用命题 \ref{prop:5.4.3}, 有 $|\alpha| = p^{k-1}$, 且点 $\bigl(\nr(\alpha/p^{k}) \otimes \nr(p^{k}/\alpha)\chi_p^{-1}\varepsilon^{k},\varepsilon^{k-1}\lambda\bigr)$ 是序的.
由命题 \ref{prop:5.3.2}, 它对应一个权 $2-k$、水平 $N$ 的非零超收敛模形式 $g$, 其 $U_p$ 特征值为 $\alpha/p^{k-1}$, Hecke 特征值系为 $\varepsilon^{k-1}\lambda$, $p$-分量特征标为 $\chi_p$.
可见 $g$ 满足陈述的条件.
\end{proof}
